\section{Monoids and automata}\label{sec:monoids}

Deck: right $M$-sets are the presheaves on $\B M$ (Example~\ref{ex:Msets}). Arches: Bridges~\ref{br:bool}, \ref{br:dm}, \ref{br:morita}.

\subsection{Excluded middle and weak excluded middle for $M$-sets}

\begin{theorem}\label{th:monoid}
Let $M$ be a monoid.
\begin{enumerate}
\item Every sub-$M$-set of every right $M$-set has a complementary sub-$M$-set iff $M$ is a group.
\item For every right $M$-set $X$, every sub-$M$-set $A$ and every $x\in X$: either $x\cdot m\notin A$ for all $m$, or for every $m$ there is $m'$ with $x\cdot mm'\in A$ --- iff $M$ satisfies the \emph{right Ore condition}: $aM\cap bM\neq\emptyset$ for all $a,b\in M$.
\end{enumerate}
\end{theorem}

\inwords{} Think of $X$ as a set of positions and of $M$ as a set of moves, $x\cdot m$ being the position reached from $x$ by the move $m$; a sub-$M$-set $A$ is a region that cannot be left. Part (1) says: the positions outside such a region also form a region that cannot be left, for every system of moves, exactly when every move can be undone. Part (2) says: from any position, either the region $A$ is unreachable, or it stays reachable whatever moves are made, exactly when any two moves can be continued so as to end at the same place ($ac = bd$).

The terminology varies: Johnstone says \q{right Ore}, and Kilp, Knauer and Mikhalev call such a monoid \emph{left reversible} \cite{hemelaer-rogers}.

\gen{} Deck $\PSh(\B M)$. Arch A: $\B M$ is a groupoid iff every element of $M$ is invertible. Arch B: a cospan in $\B M$ is a pair $a,b$, and a completing square is $c,d$ with $ac = bd$ (Example~\ref{ex:orehold}(5)). The left-hand sides are Lemma~\ref{lem:neg} and Lemma~\ref{lem:dmpointwise} unfolded for $M$-sets, where the restriction of $x$ along $m$ is $x\cdot m$.

\direct{} (1) If $M$ is a group, $X\smallsetminus A$ is a sub-$M$-set ($x\cdot m\in A$ would give $x = (x\cdot m)\cdot m^{-1}\in A$). Conversely, take $X = M$ and $A = aM$. If $a$ had no right inverse, $1\notin aM$, so $1$ would lie in the complement; the complement being closed under the action, it would contain $1\cdot a = a\in aM$, a contradiction. So every element of $M$ has a right inverse, and a monoid in which every element has a right inverse is a group (if $ab = 1$ and $bc = 1$ then $a = abc = c$, so $ba = 1$).

(2) Suppose $aM\cap bM\neq\emptyset$ always. Let $x\cdot u\in A$ and let $m\in M$. Choose $c,d$ with $uc = md$. Then $(x\cdot m)\cdot d = x\cdot md = x\cdot uc = (x\cdot u)\cdot c\in A$. Conversely, if $aM\cap bM = \emptyset$, take $X = M$, $A = aM$, $x = 1$: the move $a$ leads into $A$, but after the move $b$ no move leads into $A$, since $b\cdot m'\notin aM$.\qed

\begin{example}\label{ex:monoidlist}
\begin{enumerate}
\item \emph{Commutative monoids}, e.g.\ $(\mathbb N,+)$, $(\mathbb N_{>0},\cdot)$, $(\mathbb Z,\cdot)$, polynomials under multiplication: $ab = ba$, so Ore holds. Not groups (except trivial cases), so their topoi are De Morgan and not Boolean.
\item \emph{Monoids with a zero} $0$ ($0m = m0 = 0$): $a0 = b0$. The multiplicative monoid $M_n(k)$ of matrices, or any ring under multiplication: De Morgan.
\item \emph{Free monoids} $\Sigma^*$ with $|\Sigma|\ge 2$: Ore fails (Example~\ref{ex:orefail}). Free commutative monoids $\mathbb N^k$: Ore holds.
\item \emph{Transformation monoids.} Let $T_n$ be the maps $\{1,\dots,n\}\to\{1,\dots,n\}$ with $ab = a\circ b$. For $n\ge 2$, the constant maps $c_1, c_2$ satisfy $c_1\circ m = c_1$ for all $m$, so $c_1T_n = \{c_1\}$ and $c_2T_n = \{c_2\}$ are disjoint: Ore fails. With the opposite multiplication $a\ast b = b\circ a$, Ore holds: $a\ast c_1 = c_1\circ a = c_1 = b\ast c_1$. So the topos of right $T_n$-sets is not De Morgan, but the topos of left $T_n$-sets is. The side matters.
\item \emph{The bicyclic monoid} $B = \langle p,q\mid pq = 1\rangle$, whose elements are the words $q^ip^j$. Since $p^jq^j = 1$, we have $p^jB = B$, so $q^ip^jB = q^iB$, and $q^iB\cap q^kB\ni q^{\max(i,k)}$. Ore holds; $B$ is not a group ($qp\neq 1$).
\end{enumerate}
\end{example}

\status{} In print: (1) is \cite[Thm~2.4]{hemelaer-rogers}, (2) is part of \cite[Thm~2.31]{hemelaer-rogers}, with the commutative examples; see also Rogers' thesis \cite{rogers-thesis}. The monoids with zero are an immediate case.

\subsection{Automata}

\begin{corollary}\label{co:automata}
Let $\Sigma$ be a finite alphabet, and call a set of states of an automaton a \emph{trap} if no letter leads out of it.
\begin{enumerate}
\item If $|\Sigma| = 1$: in every automaton, for every trap $A$, every state either can never reach $A$, or can reach $A$ from every state it can reach.
\item If $|\Sigma|\ge 2$, this fails: the topos of $\Sigma$-automata is not De Morgan.
\item If $\Sigma\neq\emptyset$, the topos of $\Sigma$-automata is not Boolean.
\end{enumerate}
\end{corollary}

\begin{figure}[ht]
\centering
\begin{tikzpicture}[st/.style={circle,draw,minimum size=7mm},>=Stealth]
\node[st] (s) at (0,0) {$s$};
\node[st,fill=red!15] (a) at (2.4,0.9) {$a$};
\node[st] (b) at (2.4,-0.9) {$b$};
\draw[->] (s) -- node[above]{$x$} (a);
\draw[->] (s) -- node[below]{$y$} (b);
\draw[->] (a) edge[loop right] node{$x,y$} (a);
\draw[->] (b) edge[loop right] node{$x,y$} (b);
\node[align=left,font=\small] at (7.5,0) {trap $A = \{a\}$: from $s$ the trap is reachable (read $x$),\\ but after reading $y$ it is lost for ever.\\ So $s\notin\neg A$ and $s\notin\neg\neg A$.};
\end{tikzpicture}
\caption{The automaton that breaks the weak excluded middle over $\{x,y\}$.}\label{fig:automaton}
\end{figure}

\inwords{} For a single input letter, \q{the trap can be reached} cannot be lost by reading more input; with two letters, it can (Figure~\ref{fig:automaton}). Equivalently: over one letter, reachability of a trap is \emph{stable}; over two, it is not.

\gen{} Theorem~\ref{th:monoid} for $M = \Sigma^*$: it is commutative when $|\Sigma| = 1$, it fails Ore when $|\Sigma|\ge 2$, and it is a group only when $\Sigma = \emptyset$.

\direct{} (1) Over one letter, an automaton is a map $\tau : X\to X$ and a trap is a set with $\tau(A)\subseteq A$. If $\tau^k(x)\in A$, then from any $\tau^j(x)$, the state $\tau^{\max(j,k)}(x) = \tau^{\max(j,k)-k}(\tau^k(x))$ is reachable and lies in $A$. (2) Figure~\ref{fig:automaton}. (3) Take the automaton with two states $u\neq w$, every letter sending both to $w$; the trap $\{w\}$ has the non-trap $\{u\}$ as complement.\qed

\status{} The identification of automata with $\Sigma^*$-sets, hence the topos of automata, is standard; we have not found the De Morgan property of automata stated as such. It is an immediate case of Theorem~\ref{th:monoid}.

\begin{remark}
The same reading applies to any \q{system of commands}: a monoid $M$ of commands acting on a state space. The weak excluded middle holds for all state spaces iff any two commands $a,b$ have continuations $c,d$ that make them agree ($ac = bd$). Undo-able commands (a group) give the full excluded middle.
\end{remark}

\begin{exercise}\label{exo:automata}
Over $\Sigma = \{x,y\}$, describe all automata with at most three states for which the weak excluded middle holds for every trap. (Hint: by Lemma~\ref{lem:dmpointwise}, it is a condition on each state and each trap.)
\end{exercise}

\subsection{Transferring along a Morita equivalence of monoids}\label{ssec:transfer}

Bridge~\ref{br:morita} produces decks between monoids: $\PSh(\B M)\simeq\PSh(\B N)$ iff $\Kar(\B M)\simeq\Kar(\B N)$. Here is the basic example.

\begin{lemma}\label{lem:corner}
Let $e$ be an idempotent of $M$ such that $1\in MeM$, i.e.\ $1 = res$ for some $r,s\in M$. Then $eMe = \{eme : m\in M\}$ is a monoid with unit $e$, and right $M$-sets and right $eMe$-sets form equivalent categories.
\end{lemma}

\begin{proof}[Sketch]
$\Kar(\B(eMe))$ is the full subcategory of $\Kar(\B M)$ on the idempotents of $eMe$. It is equivalent to all of $\Kar(\B M)$, because every idempotent $f$ of $M$ is isomorphic in $\Kar(\B M)$ to an idempotent of $eMe$, namely $g = esfre$. Indeed, put $a = esf$ and $b = fre$. Then $ba = f(res)f = f$ and $ab = esfre = g$; from these, $g^2 = es f(res) fre = g$, and one checks $ga = a = af$ and $fb = b = bg$, so that $a\in gMf$ and $b\in fMg$ are inverse isomorphisms between $f$ and $g$ in $\Kar(\B M)$ (Example~\ref{ex:karM}). Bridge~\ref{br:morita} concludes.
\end{proof}

\begin{theorem}\label{th:transfer}
Let $e$ be an idempotent of $M$ with $1\in MeM$. Then $M$ satisfies the right Ore condition iff $eMe$ does, and $M$ is a group iff $eMe$ is.
\end{theorem}

\gen{} Deck: Lemma~\ref{lem:corner}. Arches: Corollary~\ref{cor:moritainv}.

\direct{} Write $1 = res$.
If $M$ satisfies Ore and $a,b\in eMe$, take $c,d$ with $ac = bd$; then $x = ecre$ and $y = edre$ lie in $eMe$ and, since $a = ae$ and $b = be$, $ax = acre = bdre = by$.
If $eMe$ satisfies Ore and $a,b\in M$, put $a' = esare$ and $b' = esbre$, in $eMe$, and take $x,y\in eMe$ with $a'x = b'y$. Multiply on the left by $r$: since $res = 1$, $a\,rex = b\,rey$.
If $M$ is a group, its only idempotent is $1$, so $e = 1$. If $eMe$ is a group and $m\in M$, let $t = ete\in eMe$ be the inverse of $g = esmre$, so $gt = e = tg$, and put $u = rets$. Using $m = (res)m$ and $e^2 = e$, $mu = r(esmre)ts = r(gt)s = res = 1$; using $m = m(res)$, $um = ret(esmre)s = r(tg)s = res = 1$. So $m$ is invertible.\qed

\begin{example}
\begin{enumerate}
\item In a \emph{finite} monoid, $1\in MeM$ forces $e = 1$: from $(re)s = 1$, the element $u = re$ has a right inverse, hence is invertible (the map $x\mapsto ux$ is surjective since $u(sx) = x$, hence injective as $M$ is finite, and $u(su) = u = u\cdot 1$ gives $su = 1$); and $ue = re^2 = u$ gives $e = 1$. So Theorem~\ref{th:transfer} is only interesting for infinite monoids.
\item In the bicyclic monoid, $e = qp$ satisfies $p\,e\,q = pqpq = 1$, and $eBe = \{q^ip^j : i,j\ge 1\}\cong B$ via $q^ip^j\mapsto q^{i-1}p^{j-1}$. (Here the theorem says nothing new, since $eBe\cong B$.)
\item Let $V$ be an infinite-dimensional vector space, $M = \End(V)$, and $e$ a projection onto a subspace $W$ with $W\cong V$. Then $1\in MeM$ and $eMe\cong\End(W)$.
\end{enumerate}
\end{example}

\status{} Morita equivalence of monoids goes back to Knauer and Banaschewski \cite{knauer,banaschewski}; in Rogers' form, $\PSh(\B M)\simeq\PSh(\B N)$ iff $N\cong eMe$ with $\beta\beta' = 1$, $\beta e = \beta$ \cite[Cor.~6.6]{rogers-1905}. The transfer of the Ore condition follows from \cite[Thm~2.31]{hemelaer-rogers} and the Morita invariance of being De Morgan; we have not found it stated for corners $eMe$.
